% !TEX root = ../enac_project_fr.tex
\chapter{Détection et traitement des sauts de cycle}
\label{ann:cycle_slips}

Cette annexe détaille la détection des sauts de cycle utilisée dans le
traitement PPP. Elle repose sur deux informations complémentaires : les
indicateurs LLI contenus dans les observations RINEX et la combinaison de
Melbourne-Wübbena calculée à partir d'observations bi-fréquences.

\section{Définition et conséquences}

Tant que le récepteur suit une porteuse sans interruption, l'ambiguïté entière
$N_{r,i}^s$ reste constante. Une perte momentanée du verrouillage peut modifier
le comptage d'un nombre entier de cycles. Entre les instants immédiatement
antérieur et postérieur à la discontinuité,

\begin{align}
N_{r,i}^s(t^+)
&=N_{r,i}^s(t^-)+\Delta N_{r,i}^s,
&
\Delta N_{r,i}^s&\in\mathbb Z\setminus\{0\},
\\
\Delta L_{r,i}^s&=\lambda_i\Delta N_{r,i}^s.
\label{ann:eq:cycle_slip_definition}
\end{align}

avec :

\begin{itemize}
    \item $t^-$ et $t^+$ Les instants situés avant et après la
    discontinuité,
    \item $\Delta N_{r,i}^s$ Le nombre entier de cycles perdus ou ajoutés,
    \item $\Delta L_{r,i}^s$ La discontinuité de phase exprimée en mètres,
    \item $\lambda_i$ La longueur d'onde de la porteuse.
\end{itemize}

La mesure de code n'est pas affectée par ce changement d'ambiguïté. En
revanche, conserver l'ancienne ambiguïté dans le filtre introduirait une erreur
systématique dans les mesures de phase. La détection est donc réalisée avant
la construction des observations de l'époque courante.

\section{Détection par LLI}

Le \textit{Loss-of-Lock Indicator} est associé aux observations de phase dans
le format RINEX~\cite{rinex304}. Il renseigne directement sur l'état du suivi
effectué par le récepteur. Les bits exploités signalent une perte de
verrouillage ou une possible ambiguïté de demi-cycle. Dans les deux cas, la
continuité de l'ambiguïté entière n'est plus garantie et le critère est

\begin{align}
\mathrm{LLI}_{r,i}^s(t_k)\neq 0
\quad\Longrightarrow\quad
\text{saut détecté}.
\label{ann:eq:lli_detection}
\end{align}

Le LLI constitue une information directe et peu coûteuse à exploiter. Il reste
toutefois dépendant du récepteur et de la manière dont celui-ci renseigne le
fichier RINEX. Il peut donc ne pas signaler certaines discontinuités, ce qui
justifie l'emploi simultané d'un indicateur construit à partir des mesures.

\section{Détection par Melbourne-Wübbena}

La combinaison de Melbourne-Wübbena utilise deux fréquences. Elle correspond à
la différence entre une combinaison de phase \textit{wide-lane} et une
combinaison de code \textit{narrow-lane}
~\cite{blewitt1990,melbourne1985,wubbena1985} :

\begin{align}
M_{\mathrm{MW},r}^s
&=
\frac{f_1L_{r,1}^s-f_2L_{r,2}^s}{f_1-f_2}
-
\frac{f_1P_{r,1}^s+f_2P_{r,2}^s}{f_1+f_2}.
\label{ann:eq:melbourne_wubbena}
\end{align}

avec :

\begin{itemize}
    \item $M_{\mathrm{MW},r}^s$ La combinaison de Melbourne-Wübbena, en
    mètres,
    \item $L_{r,1}^s$ et $L_{r,2}^s$ Les mesures de phase, en mètres,
    \item $P_{r,1}^s$ et $P_{r,2}^s$ Les mesures de code, en mètres,
    \item $f_1$ et $f_2$ Les fréquences porteuses correspondantes.
\end{itemize}

Cette combinaison élimine au premier ordre la géométrie satellite-récepteur,
les biais d'horloge, la troposphère et l'ionosphère. Elle contient donc
principalement une combinaison d'ambiguïtés, ainsi que le bruit amplifié des
mesures de code. En l'absence de saut, sa valeur demeure approximativement
constante entre deux époques consécutives. On calcule le résidu

\begin{align}
r_{\mathrm{MW},k}^s
&=M_{\mathrm{MW},r}^s(t_k)-M_{\mathrm{MW},r}^s(t_{k-1}),
\\
\left|r_{\mathrm{MW},k}^s\right|
&>\eta_{\mathrm{MW}}
\quad\Longrightarrow\quad
\text{saut détecté}.
\label{ann:eq:melbourne_wubbena_detection}
\end{align}

avec :

\begin{itemize}
    \item $r_{\mathrm{MW},k}^s$ La variation de la combinaison entre deux
    époques successives,
    \item $\eta_{\mathrm{MW}}$ Le seuil de détection.
\end{itemize}

Le seuil retenu dans ce travail est $\eta_{\mathrm{MW}}=10\,\mathrm{m}$. La
comparaison n'est effectuée que si les deux époques appartiennent au même arc
continu et si toutes les observations nécessaires sont disponibles.

\section{Fusion des indicateurs et réinitialisation}

Les critères LLI et MW sont évalués séparément pour chaque satellite. La
décision finale suit une logique inclusive : un saut est retenu dès que l'un
des deux indicateurs est déclenché. Le LLI détecte directement les pertes de
verrouillage déclarées par le récepteur, tandis que MW complète cette
information en recherchant une discontinuité dans les observations.

Lorsqu'un saut est détecté, l'ambiguïté correspondante est réinitialisée à
partir de la différence entre phase et code :

\begin{align}
\widehat N_{r,i,0}^s
&=\frac{L_{r,i}^s-P_{r,i}^s}{\lambda_i},
&
B_{r,i,0}^s
&=\lambda_i\widehat N_{r,i,0}^s.
\label{ann:eq:ambiguity_reset}
\end{align}

Sa variance est replacée à sa valeur initiale afin que le filtre puisse
réestimer rapidement ce paramètre. Cette opération reste locale à l'ambiguïté
concernée : la position, l'horloge, les paramètres troposphériques et les
ambiguïtés des autres satellites sont conservés.

Une interruption prolongée est traitée de la même manière, même en l'absence
d'un indicateur explicite :

\begin{align}
n_{\mathrm{gap}}^s
&=k-k_{\mathrm{last}}^s,
&
n_{\mathrm{gap}}^s>N_{\mathrm{out}}
&\quad\Longrightarrow\quad
\text{ambiguïté réinitialisée}.
\label{ann:eq:outage_reset}
\end{align}

avec :

\begin{itemize}
    \item $k_{\mathrm{last}}^s$ La dernière époque de présence du satellite,
    \item $n_{\mathrm{gap}}^s$ La durée de l'interruption en nombre d'époques,
    \item $N_{\mathrm{out}}$ Le seuil de réinitialisation.
\end{itemize}

La valeur retenue est $N_{\mathrm{out}}=20$ époques. Après chaque détection ou
interruption longue, les valeurs historiques utilisées par le test MW sont
également réinitialisées afin de commencer un nouvel arc de phase cohérent.
